\begin{flushright}
  \fechaclase{02 de octubre de 2026}
\end{flushright}

\subsection{Atenuación de chattering mediante observadores}

Considere el sistema no lineal de segundo orden

\begin{equation*}
    \begin{aligned}
        \dot{x}_1 &= x_2 \\
        \dot{x}_2 &= f(x_1, x_2) + u \\
        y &= x_1
    \end{aligned}
\end{equation*}

se supone que la no linealidad cumple \(|f(x,t)|\leq L\) y \(|\dot{f}(x,t)|\leq M\) para reales positivos \(L\) y \(M\). Derivando la salida se obtiene

\begin{equation*}
    \begin{aligned}
        \dot{y} &=  \dot{x}_1 = x_2 \\
        \ddot{y} &= \dot{x}_2 = f(y, t) + u\\
        \dddot{y}
    \end{aligned}
\end{equation*}

se realiza el cambio de variable

\begin{equation*}
    z_1 = y, \quad z_2 = \dot{y}, \quad z_3 = f(y, t) + u
\end{equation*}

y se reescribe el sistema como

\begin{equation*}
    \begin{aligned}
        \dot{z}_1 &= \dot{y} = z_2 \\
        \dot{z}_2 &= z_3 +bu\\
        \dot{z}_3 &= \dot{f}(x,t) \\
    \end{aligned}
\end{equation*}

o bien

\begin{equation*}
    \begin{bmatrix}
        \dot{z}_1 \\
        \dot{z}_2 \\
        \dot{z}_3
    \end{bmatrix} = \begin{bmatrix}
        0 & 1 & 0 \\
        0 & 0 & 1 \\
        0 & 0 & 0
    \end{bmatrix} \begin{bmatrix}
        z_1 \\
        z_2 \\
        z_3
    \end{bmatrix} + \begin{bmatrix}
        0 \\
        1 \\
        0
    \end{bmatrix} u + \begin{bmatrix}
        0 \\
        0 \\
        1
    \end{bmatrix} \dot{f}
\end{equation*}

\begin{equation*}
    y_z = \begin{bmatrix}
        1 & 0 & 0
    \end{bmatrix} \begin{bmatrix}
        z_1 \\
        z_2 \\
        z_3
    \end{bmatrix}
\end{equation*}

designando las matrices se obtiene

\begin{equation*}
    A=\begin{bmatrix}
        0 & 1 & 0 \\
        0 & 0 & 1 \\
        0 & 0 & 0
    \end{bmatrix}, \quad
    B=\begin{bmatrix}
        0 \\
        1 \\
        0
    \end{bmatrix}, \quad
    C=\begin{bmatrix}
        1 & 0 & 0
    \end{bmatrix} \quad
    D =\begin{bmatrix}
        0 \\ 0 \\ 1
    \end{bmatrix}
\end{equation*}

el sistema es 

\begin{equation*}
    \dot{z} = Az + Bu + D\dot{f}, \quad y_z = Cz
\end{equation*}

donde \(z_3\) es un estado extendido que representa la dinámica de la no linelaidad.\\

Se propone un observador de Luenberger tipo LESO para estimar \(f(x,t)\)

\begin{equation*}
    \begin{aligned}
        \dot{\hat{z}} &= A\hat{z} + Bu + K_oCe_o\\
        y &= Cz
    \end{aligned}
\end{equation*}

Donde el error de observación es \(e=z-\hat{z}\). Se usa la función de Lyapunov 

\begin{equation*}
    V_o=\frac{1}{2}e_o^TPe_oe_o
\end{equation*}

con la matriz \(P_o>0\). Derivando la función de Lyapunov se obtiene

\begin{equation*}
    \dot{V}_o = e_o^TP_o\dot{e}_o + \dot{e}_o^TP_o e_o 
\end{equation*}

daod que la dinámica de \(e_o\) es

\begin{equation*}
    \begin{aligned}
        \dot{e}_o &= \dot{z} - \dot{\hat{z}} = Az + Bu + D\dot{f} - A\hat{z} - Bu - K_oCe_o \\
        \dot{e}_o &= A(z-\hat{z}) + D\dot{f} - K_oCe_o \\
        \dot{e}_o &= Ae_o + D\dot{f} - K_oCe_o
    \end{aligned}
\end{equation*}

tenemos:

\begin{equation*}
    \begin{aligned}
    \dot{V}_o &= e_o^TP_o(Ae_o + D\dot{f} - K_oCe_o) + (Ae_o + D\dot{f} - K_oCe_o)^TP_oe_o \\
    \dot{V}_o &= e_o^Te_oAe_o + e_o^TA^TPe_o-e^TP_oK_oCe_o - e_o^TC^TK_oP_oe_o + e_o^TP_oD\dot{f} + (D\dot{f})^TP_oe_o \\
\end{aligned}
\end{equation*}


